\section{Methods}
\label{sec:methods}

\subsection{Framework Overview}

\methodname{} has three nested levels (\cref{fig:architecture}):
Level~1 optimizes the task; Level~1.5 adjusts search parameters;
Level~2 generates new search \emph{mechanisms} as Python code.
All levels use the same DeepSeek \texttt{deepseek-chat} model.

\begin{figure}[!htbp]
  \centering
  \includegraphics[width=\columnwidth]{architecture.pdf}
  \caption{%
    \methodname{} architecture.
    Level~1 (blue) runs the standard propose--train--evaluate loop.
    Level~1.5 (amber) adjusts search parameters every 5 iterations.
    Level~2 (green) generates new Python mechanisms via a 4-round
    research session and injects them at runtime.%
  }
  \label{fig:architecture}
\end{figure}

\subsection{Level~1: Inner Autoresearch Loop}

The inner loop implements the standard autoresearch cycle
\citep{karpathy2026autoresearch}.
At each iteration $t$:
\begin{enumerate}
  \item The LLM receives the current \texttt{train.py} (frozen at the best
        accepted configuration), the list of active editable parameters, any
        frozen parameters, and a strategic guidance string injected by
        Level~1.5.
  \item The LLM proposes a change: a set of parameter name--value pairs and a
        one-sentence hypothesis.
  \item The change is applied to a working copy of \texttt{train.py} and
        training runs for a fixed 300-second budget.
  \item If the resulting \valbpb{} is lower than the current best, the change
        is \emph{kept} (the best copy is updated); otherwise it is
        \emph{discarded}.
\end{enumerate}
The iteration budget is fixed at 30 per repeat.
The initial configuration locks \texttt{DEPTH=8} and \texttt{ASPECT\_RATIO=64}
to prevent architecture-size changes; all other parameters (\texttt{LR},
\texttt{WEIGHT\_DECAY}, \texttt{WINDOW\_PATTERN}, \texttt{HEAD\_DIM},
\texttt{TOTAL\_BATCH\_SIZE}, etc.) are editable.

\subsection{Level~1.5: Outer Search-Strategy Loop}

Level~1.5 executes every 5 inner iterations.
It receives the full trace of proposals and outcomes and produces a
\texttt{SearchConfig} update:
\begin{itemize}
  \item \textbf{Freeze} parameters that have been proposed $\geq k$ times with
        zero net improvement (default $k=3$).
  \item \textbf{Unfreeze} parameters that were frozen but have not been
        explored since the search moved to a new region.
  \item \textbf{Inject} a guidance string instructing the inner loop to
        prioritize under-explored parameters.
\end{itemize}
Level~1.5 can redirect search diversity but cannot change the proposal
generation logic, the acceptance criterion, or the loop structure.
These structural changes require Level~2.

\subsection{Level~2: Mechanism Research and Code Injection}

Level~2 executes every 2 outer cycles.
It conducts a 4-round structured dialogue, each round making a single LLM call:

\begin{enumerate}
  \item \textbf{Explore.} The LLM reads the full \texttt{runner.py} source and
        the search trace.
        It surveys mechanisms from adjacent fields (combinatorial optimization,
        online learning, design of experiments, Bayesian optimization) and
        proposes candidate improvements.
  \item \textbf{Critique.} The LLM evaluates the candidate mechanisms against
        the observed failure mode (e.g., repetitive proposals, parameter
        fixation) and selects the most promising one.
  \item \textbf{Specify.} The LLM writes a precise interface specification:
        class name, constructor arguments, key methods with signatures, and
        integration points in \texttt{runner.py}.
  \item \textbf{Generate.} The LLM writes complete, runnable Python code
        implementing the specified mechanism, including any modifications to
        \texttt{runner.py} required to call it.
\end{enumerate}

\begin{figure}[!htbp]
  \centering
  \includegraphics[width=\columnwidth]{level2_session.pdf}
  \caption{%
    Level 2 research session.
    Each session makes four LLM calls, producing a validated Python module
    that modifies the inner loop's search behavior.%
  }
  \label{fig:level2session}
\end{figure}

The generated code patches \texttt{runner.py} in place and is validated
via \texttt{importlib} dynamic loading before activation.
If the import succeeds, the patched runner replaces the active one;
if it fails, the original is restored from a pre-patch backup.
This validate-and-revert mechanism ensures that Level~2 failures are
non-destructive.

\subsection{Algorithm}

\begin{algorithm}[t]
\caption{\methodname{} (Group~C configuration)}
\label{alg:bilevel}
\begin{algorithmic}
\STATE \textbf{Input:} baseline \texttt{train.py}, runner $\phi_0$, budgets
       $T{=}30$, $K{=}5$ (outer period), $M{=}2$ (L2 period)
\STATE $\theta \leftarrow$ baseline config;\; $\phi \leftarrow \phi_0$;\;
       $t \leftarrow 0$;\; outer\_cycle $\leftarrow 0$
\WHILE{$t < T$}
  \FOR{$k = 1$ \TO $K$}
    \STATE proposal $\leftarrow$ \textsc{LLMPropose}($\theta$, $\phi$,
           guidance)
    \STATE $\theta' \leftarrow \theta \oplus$ proposal
    \STATE val $\leftarrow$ \textsc{Train}($\theta'$, budget$=$300s)
    \IF{val $<$ \textsc{BestVal}}
      \STATE $\theta \leftarrow \theta'$;\; \textsc{BestVal} $\leftarrow$ val
    \ENDIF
    \STATE $t \leftarrow t + 1$
  \ENDFOR
  \STATE guidance, frozen $\leftarrow$ \textsc{Level1.5}(trace)
  \STATE outer\_cycle $\leftarrow$ outer\_cycle $+ 1$
  \IF{outer\_cycle $\bmod M = 0$}
    \STATE $\phi' \leftarrow$ \textsc{Level2Research}($\phi$, trace)
    \IF{\textsc{ValidateImport}($\phi'$)}
      \STATE $\phi \leftarrow \phi'$
    \ELSE
      \STATE \textbf{revert} $\phi$
    \ENDIF
  \ENDIF
\ENDWHILE
\STATE \textbf{return} $\theta$
\end{algorithmic}
\end{algorithm}

\subsection{Experimental Design}

Four groups isolate the contribution of each level (\cref{tab:groups}).
All variables are held constant across groups: LLM model, GPU hardware
(RTX~5090 32\,GB, three independent servers), 300-second training budget,
30-iteration search budget, and baseline \texttt{train.py}.
Each group runs 3 independent repeats; \texttt{train.py} is restored to the
original baseline between repeats (verified by log inspection).
The primary metric is $\Delta = \text{best} - \text{baseline}$ \valbpb{}
(more negative indicates greater improvement).

\begin{table}[!htbp]
\centering
\caption{Experimental groups. All variables (LLM, GPU, budget, baseline) are held constant across groups.}
\label{tab:groups}
\begin{tabular}{@{}lll@{}}
\toprule
Group & Levels Active & Description \\
\midrule
A & Level~1 only & Pure autoresearch, no outer intervention \\
B & Level~1 + 1.5 & Inner loop plus outer strategy adjustment \\
C & Level~1 + 1.5 + 2 & Full bilevel with mechanism research \\
D & Level~1 + 2 & Inner loop plus mechanism research, no strategy adjustment \\
\bottomrule
\end{tabular}
\end{table}
